\section{Objetivo geral}
\label{sec:obj_geral}

Este trabalho tem como objetivo geral \textbf{caracterizar o viés racial em classificação facial e estudar o efeito do condicionamento por tom de pele --- escala Monk Skin Tone, com dez classes --- como sinal auxiliar arquitetural}. A pesquisa foca a disparidade de desempenho entre grupos raciais no dataset FairFace~\cite{karkkainen2021fairface} (108.501 imagens, sete classes) e investiga se o condicionamento por FiLM~\cite{perez2018film} sobre \emph{backbone} ConvNeXt-T reduz essa disparidade em comparação a \emph{baselines} públicos.

A contribuição principal projetada é \textbf{uma instância empírica documentada, conforme levantamento bibliográfico sistemático, do uso de tom de pele explícito como contexto arquitetural para \emph{race classification} multi-classe}, avaliada mediante triangulação de métricas fundamentada no Teorema da Impossibilidade~\cite{kleinberg2017inherent} e com transferência do mecanismo para verificação facial como \emph{downstream task}.

\section{Objetivos específicos}
\label{sec:obj_especificos}

Quatro objetivos específicos estruturam a execução da pesquisa em duas metas, cada um vinculado a hipóteses testáveis (Seção~\ref{sec:hipoteses}) e a contribuições esperadas (Seção~\ref{sec:contribuicoes}):

\textbf{Meta 1 — Classificador racial com fairness}

\begin{description}[font=\bfseries\sffamily]
\item[Objetivo 1.] Treinar um extrator MST próprio (ViT-B/16) sobre datasets públicos --- MSTE~\cite{monk2019scale} e Casual Conversations v2~\cite{porgali2023casual} --- e construir a matriz de distribuição cruzada Monk Skin Tone $\times$ classes raciais sobre o FairFace \emph{validation set}, com validação de concordância humana em subset estratificado. A anotação humana é realizada de forma cega e independente.

\item[Objetivo 2.] Implementar e avaliar o pipeline MST $\rightarrow$ ConvNeXt-T + FiLM $\rightarrow$ \emph{race classifier} sobre o FairFace~\cite{karkkainen2021fairface}, em comparação sistemática contra três \emph{baselines} (Seção~\ref{sec:baselines}). Inclui \emph{ablation} arquitetural com três configurações comparativas de condicionamento (Seção~\ref{sec:ablation}) sob protocolo de três sementes e \emph{bootstrap} CI.
\end{description}

\textbf{Meta 2 — Transferência para verificação}

\begin{description}[font=\bfseries\sffamily]
\item[Objetivo 3.] Transferir o encoder + FiLM treinado no Objetivo 2 para \emph{face recognition} via \emph{feature-freeze} (LAFTR,~\citeonline{madras2018learning}): congela-se o encoder e treina-se apenas uma nova cabeça de verificação sobre pares dos datasets RFW~\cite{wang2019racial} e BFW~\cite{robinson2020face}, avaliando TAR@FAR~$= 10^{-4}$ por grupo racial com \emph{threshold} fixo.

\item[Objetivo 4.] Explicitar como o condicionamento por tom de pele na tarefa de classificação (Objetivo 2) se materializa como redução de disparidade na tarefa de verificação (Objetivo 3), formalizando o mecanismo de transferência.
\end{description}

\section{Hipóteses de pesquisa}
\label{sec:hipoteses}

Quatro hipóteses testáveis, com critérios formais de confirmação e refutação, orientam o desenho experimental, descritos na Tabela \ref{tab:hipoteses}. A hipótese H4 endereça explicitamente o trabalho de \citeonline{pangelinan2023exploring}, discutido na Seção~\ref{sec:rl_fairrep}.

\begin{table}[htbp]
\centering
\small
\begin{tabular}{p{0.5cm}p{5.5cm}p{4cm}p{4cm}}
\toprule
\textbf{ID} & \textbf{Hipótese} & \textbf{Confirmação} & \textbf{Refutação} \\
\midrule
H1 & O condicionamento por tom de pele (FiLM + MST) melhora o F1 macro e reduz o \emph{Disparity Ratio} em comparação ao \emph{baseline} ResNet-34 & Melhora simultânea em ambas as métricas & Ausência de melhora em qualquer das métricas \\
H2 & O ConvNeXt-T sem condicionamento (controle arquitetural) não resolve a disparidade Latinx em relação ao ResNet-34 & Disparidade Latinx se mantém no mesmo patamar relativo & Latinx melhora substancialmente com a troca de arquitetura \\
H3 & A categoria Latinx no FairFace exibe \emph{spread} MST expressivo através de múltiplas classes da escala Monk & \emph{Spread} expressivo sobre diversas classes MST (superior ao observado em outras categorias raciais) & Spread similar ao de outras categorias, sem destaque particular \\
H4 & O encoder + FiLM treinado em classificação racial transfere propriedade de menor disparidade para verificação facial via \emph{feature-freeze} & Redução de disparidade de TAR@FAR entre grupos em RFW/BFW & Ausência de redução ou degradação após transferência \\
\bottomrule
\end{tabular}
\caption{Hipóteses de pesquisa e critérios de confirmação/refutação.}
\label{tab:hipoteses}
\end{table}

\section{Contribuições esperadas}
\label{sec:contribuicoes}

Quatro contribuições estruturais são pretendidas, organizadas nos dois eixos apresentados no Capítulo~\ref{chapter:introducao}: o \textbf{eixo fenotípico-empírico} (Contribuições 1 e 2) e o \textbf{eixo metodológico-arquitetural} (Contribuições 3 e 4).

\begin{description}[font=\bfseries\sffamily]
\item[Contribuição 1.] Extrator MST treinado sobre datasets públicos de referência (MSTE, CCv2), com protocolo de validação humana cega, aplicável a datasets faciais não anotados.

\item[Contribuição 2.] Matriz de distribuição cruzada Monk Skin Tone $\times$ classes raciais sobre o FairFace \emph{validation set}, não identificada na literatura revisada, evidenciando a variação fenotípica intra-racial.

\item[Contribuição 3.] Instância documentada na literatura revisada de FiLM-condicionamento~\cite{perez2018film} aplicado a \emph{fairness} facial em \emph{race classification} multi-classe, com \emph{ablation} arquitetural (três configurações), comparação contra três \emph{baselines} e protocolo auditável (três sementes, \emph{bootstrap} CI, triangulação de métricas).

\item[Contribuição 4.] Evidência quantitativa de transferência do condicionamento MST para verificação facial via \emph{feature-freeze}, com avaliação de disparidade de TAR@FAR por grupo racial em RFW~\cite{wang2019racial} e BFW~\cite{robinson2020face}, acompanhada de formalização do mecanismo pelo qual a melhoria no classificador se materializa na tarefa de verificação.
\end{description}